\documentclass[10pt,a4paper]{amsart}
\usepackage[margin=19mm]{geometry}
\usepackage{amsmath,amssymb,mathtools}
\usepackage[scr=rsfs,cal=euler]{mathalfa}
\usepackage{xfrac}
\usepackage[unicode,hidelinks,bookmarks=false]{hyperref}
\newcommand{\ie}{i.e.\ }
\newcommand{\cf}{cf.\ }
\newcommand{\resp}{respectively\ }
\newcommand{\Subsection}{\subsection}
\providecommand{\nobreakdash}{\nobreak\hbox{-}}
\setlength{\emergencystretch}{2em}
\multlinegap0pt
\usepackage[shortlabels]{enumitem}
\usepackage{amssymb}
\usepackage{mathtools}
\newtheorem{lemma}{Lemma}
\newtheorem{theorem}{Theorem}
\newcommand{\Wcal}{\mathcal W}
\newcommand{\sh}{\partial}
\title{Kalai's Conjecture for Tight Trees}
\author{Dhruv Mubayi \and Jacques Verstra\" ete}
\date{Source: arXiv:2609.08012v1 (CC BY 4.0)}
\begin{document}
\maketitle

\noindent\emph{Source and licence.} Kalai's Conjecture for Tight Trees. Version arXiv:2609.08012v1 (CC BY 4.0), distributed by arXiv under the Creative Commons Attribution 4.0 International licence, http://creativecommons.org/licenses/by/4.0/, as recorded in the arXiv licence metadata for this exact version and shown on the abstract page https://arxiv.org/abs/2609.08012v1. Full licence notice: \texttt{licenses/arxiv-2609.08012-CC-BY-4.0.txt}.\par\smallskip

\noindent\textit{Editorial scope.} Counted unit: the article's theorem thm:equality-shadow (source0.tex lines 696--717). The article's complete original proof of it is delivered with it (source0.tex lines 719--921, one \texttt{proof} environment of 203 lines). No mathematics is written, repaired or re-derived: the statement and the proof are the article's own lines, in the article's own notation, and no retained line is substituted or re-flowed.  Retained uncounted as the source-local context the proof needs: lines 204--209; lines 302--317; lines 634--640.  Nothing was omitted from any retained span, so the package carries no omission record.  No retained line contains a citation, so the wrapper carries no bibliography and no bibliography entry.  No retained line contains a figure environment, an image-inclusion command or drawing code, so no image is delivered and the package declares no asset dependency.  Editorial formatting only, in addition: an AMS \texttt{amsart} preamble that restates the article's own declarations and macros used by the retained lines, namely the environments \texttt{lemma}, \texttt{theorem} on separate counters, together with the standard packages those lines need.  The retained environments print in this document as Lemma 1, Theorem 1 in order of appearance, whereas the article prints the same items with the numbering of its own full document.  The complete native source of the article as distributed in the arXiv e-print for this exact version is bundled unchanged as \texttt{sources/209687/source0.tex}.
	\begin{equation}\label{eq:counts}
		\begin{aligned}
			W&=(r-1)!\,|\sh H|\,(n-r+1)!,\\
			M&=r!\,|E(H)|\,(n-r+1)!.
		\end{aligned}
	\end{equation}

	\begin{lemma}\label{lem:structure}
		Let $T$ be a tight $r$-tree with at least two edges, and fix
		$e\in E(T)$. At least one of the following holds.
		\begin{enumerate}[label=\textup{(\roman*)},leftmargin=*]
			\item\label{case:leaf}
			$T$ is obtained from a smaller tight tree $S$ by adding a new
			vertex $z$ and the edge $e=F\cup\{z\}$, where $F$ is an
			$(r-1)$-subset of an edge $f$ of $S$.
			\item\label{case:gluing}
			$T=T_1\cup T_2$ for tight trees $T_1,T_2$, each with fewer edges
			than $T$, such that
			\[
			E(T_1)\cap E(T_2)=\{e\},\qquad V(T_1)\cap V(T_2)=e.
			\]
		\end{enumerate}
	\end{lemma}

	\begin{theorem}\label{thm:rooted-count}
		For every tight $r$-tree $T$ with $t\ge1$ edges and every
		$e\in E(T)$,
		\begin{equation}\label{eq:rooted-count}
			M\le C(T,e)+(t-1)W.
		\end{equation}
	\end{theorem}

	\begin{theorem}[Equality in the shadow bound]\label{thm:equality-shadow}
		Let $r\ge2$ and $t\ge2$, let $T$ be a tight $r$-tree with $t$
		edges, and put $b=t+r-2$. If $H$ is a $T$-free $r$-graph, then
		\begin{equation}\label{eq:equality-shadow}
			r|E(H)|=(t-1)|\sh H|
		\end{equation}
		if and only if the applicable one of the following conditions holds.
		\begin{enumerate}[label=\textup{(\roman*)},leftmargin=*]
			\item If $T$ is a sunflower, then
			\[
			d_H(A)=t-1\qquad\text{for every }A\in\sh H.
			\]
			\item If $T$ is not a sunflower, then there is a family
			$\mathcal B\subseteq\binom{V(H)}b$ such that
			\begin{equation}\label{eq:equality-blocks}
				E(H)=\bigcup_{B\in\mathcal B}\binom Br,
				\qquad |B\cap B'|\le r-2
				\quad\text{for distinct }B,B'\in\mathcal B.
			\end{equation}
		\end{enumerate}
		Vertices outside the union of the blocks are allowed and are isolated.
	\end{theorem}

	\begin{proof}
		The assertion is immediate when $E(H)=\varnothing$, with
		$\mathcal B=\varnothing$ in the second case. Assume henceforth
		that $E(H)\ne\varnothing$ and that~\eqref{eq:equality-shadow}
		holds. With the notation of~\eqref{eq:counts}, equality says
		\begin{equation}\label{eq:equality-MW}
			M=(t-1)W.
		\end{equation}
		
		\medskip\noindent
		\emph{Equality forces constant positive codegree.}
		Fix a tight-tree construction of $T$, and write its last edge as
		$e=F\cup\{\ell\}$, where $\ell$ is the newly introduced vertex.
		Delete $e$ and $\ell$ to obtain a tight tree $S$ with $t-1$ edges,
		and choose an edge $f=F\cup\{u\}$ of $S$. Order $f$ with the
		vertices of $F$ first and $u$ last. Theorem~\ref{thm:rooted-count}
		and~\eqref{eq:equality-MW} give
		\begin{equation}\label{eq:equality-rooted-lower}
			C(S,f)\ge M-(t-2)W=W.
		\end{equation}
		
		For each $\pi\in\Wcal$, there is at most one good state for
		$(S,f)$. Indeed, if $i<j$ are two good positions, choose an
		embedding witnessing the state at $i$. Its image lies in the
		first $i$ entries, and it maps $F$ to the first $r-1$ entries.
		The vertex $v_j$ is unused and completes these entries to an
		edge of $H$. Sending $\ell$ to $v_j$ therefore extends the
		embedding to $T$, a contradiction. Together with
		\eqref{eq:equality-rooted-lower}, this proves that every
		$\pi\in\Wcal$ has exactly one good state for $(S,f)$.
		
		\bigskip
		
		Suppose that $A\in\sh H$ satisfies $d_H(A)\le t-2$. Choose
		a permutation beginning with the vertices of $A$, followed by
		all vertices of $N_H(A)$, followed by the remaining vertices.
		Every marked position $j$ in this permutation satisfies
		\[
		j\le r-1+d_H(A)\le t+r-3.
		\]
		Since $S$ has $t+r-2$ vertices, none of these positions can be
		good for $(S,f)$. This is a contradiction. Hence every member
		of $\sh H$ has codegree at least $t-1$. Double counting gives
		\[
		\sum_{A\in\sh H}d_H(A)=r|E(H)|=(t-1)|\sh H|,
		\]
		so in fact
		\begin{equation}\label{eq:equality-codegree}
			d_H(A)=t-1\qquad\text{for every }A\in\sh H.
		\end{equation}
		
		If $T$ is a sunflower, this proves necessity. Conversely,
		\eqref{eq:equality-codegree} forbids that sunflower and gives
		\eqref{eq:equality-shadow} by double counting. This proves
		the first case of the theorem.
		
		\medskip\noindent
		\emph{For a nonsunflower, every codimension-two link is a union
			of cliques.}
		Assume now that $T$ is not a sunflower. Use the construction
		already chosen to form the parent tree $D$ on $E(T)$: join each
		edge other than the first to its chosen parent. Adjacent nodes
		meet in exactly $r-1$ vertices, and for each vertex $v$ of $T$
		the nodes containing $v$ form a connected subtree. These are
		the properties established in the proof of
		Lemma~\ref{lem:structure}. The node $e$ is a leaf, and $D-e$
		represents $S$.
		
		\bigskip
		
		The nodes of $D-e$ containing all of $F$ form a nonempty
		connected subtree: they are the intersection of the subtrees
		corresponding to the vertices of $F$. This subtree is proper,
		since otherwise every edge of $T$ would contain $F$ and $T$
		would be a sunflower. There are therefore adjacent nodes
		$f_0,g_0$ in $D-e$ with $F\subseteq f_0$ and
		$F\nsubseteq g_0$. They have the form
		\begin{equation}\label{eq:equality-two-edges}
			f_0=F\cup\{u_0\},
			\qquad
			g_0=(F\setminus\{a\})\cup\{u_0,w\},
		\end{equation}
		where $a\in F$ and $w\notin f_0$.
		Rooting $D-e$ at $f_0$ and placing $g_0$ second gives a
		tight-tree construction of $S$ beginning with $f_0,g_0$.
		Indeed, in any order with parents preceding children, the
		vertex of a child edge outside its parent cannot have occurred
		earlier: otherwise the connectedness of its occurrence nodes
		would force it to belong to the parent.
		
		\bigskip
		
		For an $(r-2)$-set $Q\subseteq V(H)$, let $L_H(Q)$ be the
		graph on $V(H)\setminus Q$ with
		\[
		xy\in E(L_H(Q))
		\quad\Longleftrightarrow\quad
		Q\cup\{x,y\}\in E(H).
		\]
		Suppose that $x,y,z$ induce a three-vertex path in $L_H(Q)$.
		Thus
		\begin{equation}\label{eq:equality-link-path}
			Q\cup\{x,y\},\ Q\cup\{y,z\}\in E(H),
			\qquad Q\cup\{x,z\}\notin E(H).
		\end{equation}
		Embed the two edges in~\eqref{eq:equality-two-edges} by an
		arbitrary bijection from $F\setminus\{a\}$ to $Q$ and by
		\[
		a\mapsto x,\qquad u_0\mapsto y,\qquad w\mapsto z.
		\]
		Extend this embedding greedily along the construction of $S$
		just described. After $q$ edges have been embedded, there are
		$q+r-1$ used vertices. The attaching $(r-1)$-face of the next
		edge belongs to an already embedded edge, so its image lies in
		$\sh H$ and has codegree $t-1$ by~\eqref{eq:equality-codegree}.
		Exactly $q\le t-2$ used vertices lie outside that face. An
		unused extension therefore exists at every step.
		
		\bigskip
		
		Let $\varphi$ be the resulting embedding of $S$, and set
		$A=\varphi(F)=Q\cup\{x\}$. The image of $S$ contains exactly
		$t-1$ vertices outside $A$. One of them is $z=\varphi(w)$,
		which is not in $N_H(A)$ by~\eqref{eq:equality-link-path}.
		Thus at most $t-2$ vertices of $N_H(A)$ are used. Since
		$d_H(A)=t-1$, an unused neighbour of $A$ remains; sending
		$\ell$ to that vertex produces a copy of $T$. This is a
		contradiction.
		
		\bigskip
		
		It follows that every $L_H(Q)$ contains no induced
		three-vertex path. Each of its connected components is
		therefore complete: a shortest path between two nonadjacent
		vertices would otherwise begin with an induced three-vertex
		path. Moreover, every vertex of positive degree in $L_H(Q)$
		has degree $t-1$, since that degree is
		$d_H(Q\cup\{x\})$ for the corresponding vertex $x$.
		Consequently every nontrivial component of every $L_H(Q)$
		is a copy of $K_t$.
		
		\medskip\noindent
		\emph{The local cliques determine complete blocks.}
		For $A\in\sh H$, put
		\[
		B_A=A\cup N_H(A).
		\]
		By~\eqref{eq:equality-codegree}, $|B_A|=t+r-2=b$.
		Choose $a\in A$ and $x\in N_H(A)$, and let
		\[
		Q=A\setminus\{a\},\qquad A'=Q\cup\{x\}.
		\]
		The component of $L_H(Q)$ containing $a$ is a clique with
		vertex set $\{a\}\cup N_H(A)$. The vertex $x$ belongs to
		this clique, so
		\[
		N_H(A')=
		\bigl(\{a\}\cup N_H(A)\bigr)\setminus\{x\}.
		\]
		In particular, $A'\in\sh H$ and
		\begin{equation}\label{eq:equality-block-exchange}
			B_{A'}=B_A.
		\end{equation}
		
		Every $(r-1)$-subset of $B_A$ can be reached from $A$ by
		successive single-vertex exchanges. At each step, the entering
		vertex belongs to the current block but not to the current
		face, and hence belongs to the neighbourhood of that face.
		Repeated application of~\eqref{eq:equality-block-exchange}
		therefore shows that
		\begin{equation}\label{eq:equality-all-faces}
			A'\in\sh H,\qquad B_{A'}=B_A
			\quad\text{for every }A'\in\binom{B_A}{r-1}.
		\end{equation}
		It follows that every $r$-subset of $B_A$ is an edge of $H$.
		Conversely, every edge of $H$ is contained in $B_A$ for each
		of its $(r-1)$-faces $A$.
		
		Let $\mathcal B$ be the family of distinct sets $B_A$ as
		$A$ ranges over $\sh H$. The preceding paragraph proves
		\[
		E(H)=\bigcup_{B\in\mathcal B}\binom Br.
		\]
		If two members $B,B'$ shared an $(r-1)$-set $A$, then
		\eqref{eq:equality-all-faces} would give $B=B_A=B'$.
		Thus distinct blocks intersect in at most $r-2$ vertices,
		as required in~\eqref{eq:equality-blocks}.
		
		\medskip\noindent
		\emph{Conversely, the block construction attains equality.}
		Suppose $H$ has the form~\eqref{eq:equality-blocks}. Each
		edge belongs to a unique block, and edges belonging to
		different blocks cannot meet in $r-1$ vertices. Following
		the parent edges in a tight-tree construction, any copy of
		$T$ in $H$ would therefore lie in a single block. This is
		impossible, since each block has $t+r-2$ vertices and $T$
		has $t+r-1$ vertices.
		
		Each member of $\sh H$ belongs to exactly one block, and
		its neighbourhood consists of the $b-(r-1)=t-1$ remaining
		vertices of that block. Thus~\eqref{eq:equality-codegree}
		holds, and double counting proves~\eqref{eq:equality-shadow}.
	\end{proof}

\end{document}
